\begin{definition}\label{def:LL}
For $f\ge 1$, define $q_{-}=\lfloor\frac{f}{2}\rfloor$ and $q^{+}=\lceil\frac{f}{2}\rceil$;
note that $q_{-}+q^{+}=f$ for every $f$ and $q_{-}=q^{+}$ for even $f$.

Then the function  $\LL_{f}$ on
$d\geq 1$ is defined as follows: 
\begin{itemize}
\item[(1)] $\LL_f(d)=f$ if $1\leq d\leq f-1$;
\end{itemize}
and for $d\geq f$, blocks of length $2f$ will be defined,
for every $k\geq 1$ by:
\begin{itemize}
\item[(2)] $\LL_f(d)=k(f+1+j)$, if $d=(2k-1)f+j$ for $0\leq j\leq q_{-}$;
\item[(3)] if $f$ is odd then $\LL_f((2k-1)f+q^{+})=q^{+}$;
\item[(4)] 
$\LL_f(d)=k(f+1+q_{-})+q^{+}+1-j\cdot (k-1)$, if $d=(2k-1)f+q^{+}+1+j$ for $0\leq j\leq q_{-}+1$;
\item[] with the correction that, if $f\leq 4$ or $k>1$ then $k$ 
is to be subtracted from $\LL_f(2kf-1)$;
\item[(5)] $\LL_f(d)=(k+1)(f-1)+3$ if $2kf+3\leq d\leq (2k+1)f-1$;
\item[] with the correction for $k=1$ that 1 is to be added to $\LL_f(d)$, and
\item[] with the correction (for $k=2$) that $\LL_f(4f+3)=3f+1$ instead of $3f$.
\end{itemize}
To this general definition of $\LL_f$ the following
modifications are made for small values of the parameters:
\begin{itemize}
\item for $f=1$:
\begin{itemize}
\item[] $\LL_1(3)=4$;
\item[] add $\frac{d+7}{2}$ to $\LL_1(d)$ for $d\equiv 1\bmod 2$ with $d\geq 7$;
\end{itemize}
\item for $f=2$:
\begin{itemize}
\item[] $\LL_2(2)=1$;
\item[] $\LL_2(6)=2$;
\item[] add $\frac{d+5}{4}$ to $\LL_2(d)$ for $d\equiv 3\bmod 4$;
\item[] add 1 to $\LL_2(d)$ for $d\equiv 0\bmod 4$.
\end{itemize}
\item for $f=3$:
\begin{itemize}
\item[] $\LL_3(d)=2$ for $d\equiv 5\bmod 6$;
\item[] $\LL_3(6)=4$;
\item[] $\LL_3(7)=2$;
\item[] $\LL_3(8)=4$;
\item[] $\LL_3(9)=4$;
\item[] $\LL_3(10)=5$;
\item[] $\LL_3(15)=26$;
\item[] add $\frac{d+20}{6}$ to $\LL_3(d)$ for $d\equiv 4\bmod 6$ 
with $d\geq 16$;
\item[] add 4 to $\LL_3(d)$ for $d\equiv 3\bmod 6$ 
with $d\geq 21$.
\end{itemize}
\item for $f=4$:
\begin{itemize}
\item[] $\LL_4(5)=2$;
\item[] $\LL_4(36)=5$;
\item[] add $\frac{d+44}{8}$ to $\LL_4(d)$ for $d\equiv 4\bmod 8$ 
with $d\geq 44$.
\end{itemize}
\end{itemize}
\end{definition}
\begin{definition}\label{def:KK}
Define, for the residue classes $r\equiv d\bmod 2f$, 
fixing the additional notation
$d=r+h\cdot (2f)$, and choosing representatives $1\leq r\leq 2f$,
a function $\KK_{r}$ of $h$ and $f$.

\begin{itemize}
\item[] For $f\geq 5$
the function is 
\begin{itemize}
\item[] linear in $h$ for fixed $f$ for the following $f+3$ residue classes $r\bmod 2f$:
\begin{itemize}
\item[(a)] $\KK_{1}(h,f)=(2f+1)\cdot h+(2f+7)$;
\item[(b)] $\KK_{2}(h,f)=(2f)\cdot h+(2f+7)$;
\item[(c)] $\KK_{r}(h,f)=(f+1)\cdot h+(3f+1+r)$
for $r=3, 4,\ldots,f-1$;
\item[(d)] $\KK_{f}(h,f)=(4f+1)\cdot h+(5f-3)$;
\item[(e)] $\KK_{f+1}(h,f)=(2f+3)\cdot h+(2f+7)$;
\item[(f)] $\KK_{2f-1}(h,f)=f\cdot h+(2f+1)$;
\item[(g)] $\KK_{2f}(h,f)=f\cdot h+(f+2)$.
\end{itemize}
\item[] and quadratic in $h$ for fixed $f$ for the remaining $f-3$ cases, as given below,
using the notation $q=\lfloor\frac{f+1}{2}\rfloor$ and $s=r-(f+q)$:
\begin{itemize}
\item[(h)] for even $f\geq 6$, and $f+3\leq r\leq 2f-2$:
\begin{itemize}
\item[(h$_0$)] $\KK_{f+q}(h,f)=\frac{9f+14}{4}h^2+\frac{15f+42}{4}h+
\frac{3f+2}{2}$;
\item[(h$_1$)] $\KK_{r}(h,f)=\frac{9f+16-2s}{4}h^2+\frac{15f+40+10s}{4}h+\frac{3f+2s}{2}$ when $s>0$; %$r>f+q$;
\item[(h$_2$)] $\KK_{r}(h,f)=\frac{9f+16+2s}{4}h^2+
\frac{15f+42+6s}{4}h+\frac{3f+2+2s}{2}$
when $s<0$; %$r<f+q$;
\end{itemize}
\item[(i)] while for odd $f\geq 5$, and $f+3\leq r\leq 2f-2$:
\begin{itemize}
\item[(i$_0$)] $\KK_{f+q}(h,f)=\frac{9f+15}{4}h^2+\frac{15f+45}{4}h+
\frac{3f-1}{2}$;
\item[(i$_1$)] $\KK_{r}(h,f)=\frac{9f+15-2s}{4}h^2+
\frac{15f+45+10s}{4}h+\frac{3f+1+2s}{2}$
when $s>0$; %$r>f+q$;
\item[(i$_2$)] $\KK_{r}(h,f)=\frac{9f+15+2s}{4}h^2+
\frac{15f+45+6s}{4}h+\frac{3f+3+2s}{2}$
when $s<0$; %$r<f+q$;
\end{itemize}
\item[(j)] when $r=f+2$ alternatingly, depending on $d\bmod 4f$:
\begin{itemize}
\item[(j$_0$)] $\KK_{r}(h,f)=(4f+9)h^2+(4f+9)h$
when $d\equiv f+2\bmod 4f$;
\item[(j$_1$)] $\KK_{r}(h,f)=(4f+9)h^2+(4f+9)h-1$
when $d\equiv 2f+f+2\bmod 4f$.
\end{itemize}
\item[(k)] Finally, for small values of $d$ (that is, of $h$ and $r$),
the following corrections are to be made to $\KK_{r}$ for
$f\geq 5$:
\begin{itemize}
\item[(k$_0$)] if $h=0$, subtract $2f+r$ from $\KK_{r}$;
\item[(k$_1$)] if $h=1$, subtract $3f+r$ from $\KK_{r}$;
\item[(k$_2$)] if $r=1$ and $h=0$, add $f-5$ to $\KK_{1}$;
\item[(k$_3$)] if $r=1$ and $h=1$, subtract $5$ from $\KK_{1}$;
\item[(k$_4$)] if $r=1$ and $h=2$ subtract $4f+7$ from $\KK_{1}$;
\item[(k$_5$)] if $r=2$ and $h=0$ add $f-4$ to $\KK_{2}$;
\item[(k$_6$)] if $r=2$ and $h=1$, subtract $3$ from $\KK_{2}$;
\item[(k$_7$)] if $r=3$ and $h=2$, add $4f+1$ to $\KK_{3}$.
\end{itemize}
\end{itemize}
\end{itemize}
\item[] For $1\leq f\leq 4$ we define
\begin{itemize}
\item[$\bullet$] for $f=1$:
\begin{itemize}
\item[] $$\KK_{1}(1+2h,1)=
\begin{cases}
2,&\text{if\ } h=0;\\
4,&\text{if\ } h=1;\\
5,&\text{if\ } h=2;\\
h+4,&\text{if\ } h\geq 3.
\end{cases}$$
\item[] $$\KK_{2}(2+2h,1)=
\begin{cases}
0,&\text{if\ } h=0;\\
1,&\text{if\ } h\geq 1.
\end{cases}$$
\end{itemize}
\item[$\bullet$] for $f=2$:
\begin{itemize}
\item[] $$\KK_{1}(1+4h,2)=
\begin{cases}
3,&\text{if\ } h=0;\\
4,&\text{if\ } h=1;\\
6,&\text{if\ } h=2;\\
5h+11,&\text{if\ } h\geq 3.
\end{cases}$$
\item[] $$\KK_{2}(2+4h,2)=
\begin{cases}
2,&\text{if\ } h=0;\\
5,&\text{if\ } h=1;\\
7,&\text{if\ } h=2;\\
5h+11,&\text{if\ } h\geq 0.
\end{cases}$$
\item[] $$\KK_{3}(3+4h,2)=2h+5.$$
\item[] $$\KK_{4}(4+4h,2)=2h+6.$$
\end{itemize}
\item[$\bullet$] for $f=3$:
\begin{itemize}
\item[] $$\KK_{1}(1+6h,3)=
\begin{cases}
4,&\text{if\ } h=0;\\
5,&\text{if\ } h=1;\\
8,&\text{if\ } h=2;\\
7h+13,&\text{if\ } h\geq 3.
\end{cases}$$
\item[] $$\KK_{2}(2+6h,3)=
\begin{cases}
4,&\text{if\ } h=0;\\
5,&\text{if\ } h=1;\\
6h+13,&\text{if\ } h\geq 2.
\end{cases}$$
\item[] $$\KK_{3}(3+6h,3)=
\begin{cases}
4,&\text{if\ } h=0;\\
19,&\text{if\ } h=1;\\
36,&\text{if\ } h=2;\\
81,&\text{if\ } h=3;\\
10h+19,&\text{if\ } h\geq 4.
\end{cases}$$
\item[] $$\KK_{4}(4+6h,3)=
\begin{cases}
4,&\text{if\ } h=0;\\
9h+14,&\text{if\ } h\geq 1.
\end{cases}$$
\item[] $$\KK_{5}(5+6h,3)=
\begin{cases}
9,&\text{if\ } h=0;\\
26,&\text{if\ } h=1;\\
\frac{21}{2}h^2+\frac{27}{2}h,&\text{if\ } h\geq 2.
\end{cases}$$
\item[] $$\KK_{6}(6+6h,3)=3h+5.$$
\end{itemize}
\item[$\bullet$] for $f=4$:
\begin{itemize}
\item[] $$\KK_{1}(1+8h,4)=
\begin{cases}
5,&\text{if\ } h=0;\\
6,&\text{if\ } h=1;\\
10,&\text{if\ } h=2;\\
9h+15,&\text{if\ } h\geq 3.
\end{cases}$$
\item[] $$\KK_{2}(2+8h,4)=
\begin{cases}
5,&\text{if\ } h=0;\\
6,&\text{if\ } h=1;\\
8h+15,&\text{if\ } h\geq 2.
\end{cases}$$
\item[] $$\KK_{3}(3+8h,4)=
\begin{cases}
5,&\text{if\ } h=0;\\
6,&\text{if\ } h=1;\\
43,&\text{if\ } h=2;\\
5h+16,&\text{if\ } h\geq 3.
\end{cases}$$
\item[] $$\KK_{4}(4+8h,4)=
\begin{cases}
5,&\text{if\ } h=0;\\
12,&\text{if\ } h=1;\\
116,&\text{if\ } h=2;\\
112,&\text{if\ } h=3;\\
193,&\text{if\ } h=4;\\
32h+40,&\text{if\ } h\geq 5.
\end{cases}$$
\item[] $$\KK_{5}(5+8h,4)=
\begin{cases}
5,&\text{if\ } h=0;\\
13,&\text{if\ } h=1;\\
11h+15,&\text{if\ } h\geq 2.
\end{cases}$$
\item[] $$\KK_{6}(6+8h,4)=
\begin{cases}
5,&\text{if\ } h=0;\\
14,&\text{if\ } h=1;\\
20,&\text{if\ } h=2;\\
51,&\text{if\ } h=3;\\
\frac{25}{4}h^2-7,&\text{if\ } h\geq 4,\text{\ even},\\
\frac{25}{4}h^2-\frac{29}{4}&\text{if\ } h\geq 5\text{\ odd}.
\end{cases}$$
\item[] $$\KK_{7}(7+8h,4)=
\begin{cases}
11,&\text{if\ } h=0;\\
4h+9,&\text{if\ } h\geq 1.
\end{cases}$$
\item[] $$\KK_{8}(8+8h,4)=4h+6.$$
\end{itemize}
\end{itemize}
\end{itemize}
\end{definition}
